% !TeX root = ../portafoliofinal.tex
% !TeX encoding = utf8

\chapter{Resumen ejecutivo}

El presente portafolio grupal final constituye el trabajo recopilatorio de la
asignatura \emph{Dirección de Recursos Humanos~I} (Doble Grado ADE-Ingenierías,
Universidad de Granada, curso 2025--2026). Recoge, de forma integrada y enriquecida,
las tres prácticas previas sobre la empresa \textbf{CaixaBank, S.A.}, junto con
un informe global de evaluación y una autoevaluación del equipo.

\medskip

\textbf{Temas 2 y 3 — Análisis, Diseño y Planificación de RRHH.}
Se analiza cómo CaixaBank realiza el Análisis del Puesto de Trabajo (APT) mediante
la Secretaría Técnica de Negocios y la consultora Deloitte, con un enfoque
\textit{bottom-up} orientado a competencias. Se describen tres puestos representativos
(Director/a de Oficina, Gestor/a de Banca Privada y Asesor/a Comercial) y se estudian
las formas de rediseño aplicadas, en especial la rotación interna y el teletrabajo.
En planificación de RRHH, se examina la gestión estratégica del dimensionamiento
apoyada en SAP SuccessFactors y Microsoft Copilot, el desequilibrio de oferta
en perfiles tecnológicos, los criterios del SIRH y los procesos sustractivos (ERE
2021 con Bankia, prejubilaciones y \textit{outplacement} parcial).

\medskip

\textbf{Temas 4 y 5 — Reclutamiento, Selección e Integración.}
Se detalla el modelo mixto de reclutamiento de CaixaBank, que combina la movilidad
interna (portal de vacantes) con canales externos (LinkedIn, InfoJobs, ETTs,
\textit{CaixaBank Talent Program}). Se analiza el ciclo completo de selección
—preselección ATS, pruebas diferenciadas por perfil, dinámicas de grupo,
\textit{assessment centers} y videoentrevistas— así como el proceso de socialización
formal e informal y el manual de acogida. Se incorporan propuestas de mejora basadas
en \textit{employer branding} segmentado, gobernanza ética de la IA y gamificación.

\medskip

\textbf{Temas 6, 7 y 8 — Formación, Evaluación del Rendimiento y Retribución.}
Se examina el ciclo formativo completo de la entidad: identificación de necesidades
por niveles, diseño mixto (interno/externo), implantación en la plataforma
\textit{Evoluciona} y evaluación mediante el modelo Kirkpatrick. Se estudian
las tendencias adoptadas: \textit{e-learning}, modelo 70-20-10, universidades
corporativas, \textit{outdoor training} y gamificación. En gestión de carrera se
describen las fases de valoración, dirección y desarrollo, con especial atención
al \textit{coaching} ejecutivo (método GROW), \textit{mentoring} y \textit{mentoring}
inverso. En evaluación del rendimiento se analizan métodos objetivos, subjetivos y
mixtos, las estrategias de reforzamiento y las tendencias de evaluación continua y
por compromisos ASG. Finalmente, en retribución se describe la estructura salarial
de cinco puestos conforme al IV Convenio Colectivo de Banca, los tres tipos de
equidad, el modelo híbrido puesto/competencias, los incentivos y las políticas
de retribución flexible y salario emocional.

\medskip

\textbf{Informe global.}
Se presentan recomendaciones sistémicas que interrelacionan las distintas prácticas
de recursos humanos analizadas, apostando por un modelo integrado y sostenible de
gestión del talento en CaixaBank. Se incluyen las limitaciones de información
identificadas durante el proceso de investigación.

\medskip

\textbf{Palabras clave:} análisis de puestos, planificación de RRHH, reclutamiento,
selección, integración, formación corporativa, \textit{e-learning},
gestión de carrera, evaluación del rendimiento, retribución flexible,
salario emocional, CaixaBank, convenio colectivo de banca.

\cleardoublepage
\endinput
